% =============================================================================
\chapter{Wprowadzenie}
\label{ch:introduction}
% =============================================================================

Duże modele językowe (LLM) stały się codziennym narzędziem programistów. Uzupełniają kod
w edytorach, odpowiadają na pytania o biblioteki i coraz częściej piszą całe funkcje lub moduły na
podstawie krótkiego opisu. Ocena tych umiejętności podążała tą samą drogą. Wczesne benchmarki, takie
jak HumanEval~\cite{chen2021codex} i MBPP~\cite{austin2021mbpp}, mierzą, jak często model tworzy
funkcję w Pythonie, która przechodzi kilka testów jednostkowych. Późniejsze prace przetłumaczyły te
problemy na inne języki~\cite{cassano2023multiple,zheng2023codegeex,athiwaratkun2023mbxp}, dodały
mocniejsze testy~\cite{liu2023evalplus} lub przeszły do zadań na poziomie całych
repozytoriów~\cite{jimenez2024swebench}.

Niemal wszystkie te prace traktują język programowania jako szczegół powierzchniowy: zadanie
pozostaje takie samo, zmienia się tylko składnia. Języki różnią się jednak nie tylko składnią.
\emph{Paradygmat programowania}~\cite{floyd1979paradigms,vanroy2004ctm}, czy to imperatywny,
obiektowy, funkcyjny czy logiczny, określa, jak rozkłada się rozwiązanie na części, jak wyraża się
stan i przepływ danych, jak reprezentuje się brak wartości i błędy oraz kto zarządza pamięcią.
Algorytm Dijkstry napisany z modyfikowalnym kopcem binarnym w C, z niemodyfikowalnymi mapami
i zwinięciami (ang.~\emph{folds}) w Haskellu albo z relacjami i nawracaniem w Prologu to trzy różne
ćwiczenia. Dane treningowe modeli LLM są ponadto silnie przesunięte w stronę kilku popularnych
języków~\cite{kocetkov2022stack}, a modele radzą sobie wyraźnie gorzej z językami o małych zasobach
danych~\cite{cassano2023multiple,cassano2024knowledge}. Smalltalk, Pascal i Prolog są takimi
językami pod każdym względem.

W niniejszym raporcie pytamy, czy te różnice mają znaczenie dla kodu generowanego przez LLM.
Utrzymujemy stałe \emph{zadanie} i \emph{testy}, a zmieniamy tylko język, a wraz z nim paradygmat.
Następnie mierzymy poprawność funkcjonalną, zdolność modelu do naprawy własnych błędów na podstawie
informacji zwrotnej z kompilatora i testów oraz idiomatyczność kodu końcowego. Idiomatyczność jest
istotna, ponieważ rozwiązanie może przejść wszystkie testy, a mimo to być „C napisanym
w Haskellu”, czyli kodem, który programista utrzymujący kod w tym języku by odrzucił.

\section{Pytania badawcze}
\label{sec:rqs}

\begin{rqbox}
\begin{description}[leftmargin=3.2em,labelwidth=2.8em,itemsep=3pt]
  \item[\RQ{1}] \textbf{Poprawność.} Jak różni się poprawność funkcjonalna kodu generowanego przez
    LLM za pierwszym podejściem między językami i paradygmatami, gdy zadanie i testy pozostają
    stałe?
  \item[\RQ{2}] \textbf{Samonaprawa.} W jakim stopniu luka ta się zmniejsza, gdy model otrzymuje
    komunikaty kompilatora, logi testów i nazwy niezaliczonych testów, i czy ta zdolność zależy od
    paradygmatu?
  \item[\RQ{3}] \textbf{Idiomatyczność.} Czy wygenerowany kod jest idiomatyczny dla docelowego
    języka i paradygmatu w ocenie zaślepionego sędziego LLM posługującego się stałą rubryką?
  \item[\RQ{4}] \textbf{Praktyczność.} Ile kosztuje taka ewaluacja i jakie przeszkody praktyczne
    (wykonywanie instrukcji, format odpowiedzi, sposób hostowania modelu) wpływają na jej
    trafność?
\end{description}
\end{rqbox}

\section{Wkład}

\begin{enumerate}
  \item \textbf{Benchmark obejmujący wiele paradygmatów.} Dziesięć zadań algorytmicznych, każde
    z dwoma algorytmami i 20 testami jednostkowymi, zaimplementowanych w siedmiu językach
    reprezentujących cztery paradygmaty (łącznie 1\,400 przypadków testowych). Każde zadanie zawiera
    specyfikację, kontrakt, szkielet (ang.~\emph{stub}) i ukryty zestaw testów. Każdy język ma
    własny obraz Docker, a błędy pamięci są sprawdzane Valgrindem tam, gdzie język na to pozwala
    (\cref{ch:benchmark}).
  \item \textbf{System ewaluacyjny.} Program w Pythonie, który generuje prompty z szablonów dla
    poszczególnych języków, wywołuje model, ocenia odpowiedź w kontenerach, klasyfikuje wynik do
    jednej z dziewięciu kategorii, prowadzi pętlę samonaprawy, ocenia idiomatyczność za pomocą
    zaślepionego sędziego LLM i agreguje raport. Przebiegi można wznawiać, są w pełni śledzone,
    a ich koszt jest kontrolowany (\cref{ch:harness}).
  \item \textbf{Badanie empiryczne.} Pełny przebieg modelu \haiku{} na wszystkich 70 parach
    język--zadanie, oceniony przez \sonnet{}, wraz z analizą poprawności, rodzajów błędów,
    samonaprawy, idiomatyczności i kosztów (\cref{ch:results,ch:discussion}).
  \item \textbf{Wyniki negatywne i wnioski.} Zadania w stylu HumanEval okazały się zbyt łatwe, by
    rozróżnić paradygmaty, a lokalne modele o otwartych wagach były niepraktyczne ze względu na
    szybkość i format odpowiedzi (\cref{ch:preliminary}). Przebieg ujawnił też dwa defekty
    w zestawie testów dla Prologu (\cref{sec:benchmark-defects}).
\end{enumerate}

\section{Struktura raportu}

\Cref{ch:background} przedstawia paradygmaty oraz prace pokrewne dotyczące benchmarków generowania
kodu, samonaprawy i sędziów LLM. \Cref{ch:preliminary} opisuje nasze eksperymenty pilotażowe
i wyciągnięte z nich wnioski. \Cref{ch:benchmark} prezentuje benchmark, a \cref{ch:harness} system
ewaluacyjny, w tym konstrukcję promptów. \Cref{ch:results} zawiera wyniki, \cref{ch:discussion}
ich omówienie wraz z zagrożeniami dla trafności, a \cref{ch:conclusions} podsumowanie. Dodatki
zawierają pełne prompty, rubrykę sędziego i wyniki dla poszczególnych zadań.
